% ch03-electronics.tex -- Chapter 3: Series, parallel and Kirchhoff's laws
\chapter{Series, parallel and Kirchhoff's laws}
\label{ch:networks}

Chapter~2 gave you one rule, $V=IR$, and one tool, ngspice. This chapter
connects resistors into \emph{networks} and asks the obvious question: what is
the voltage \emph{here}, when the circuit has more than one place to look?
You will learn the two bookkeeping laws every engineer uses without thinking
(Kirchhoff's), the voltage divider — the single most-used fragment of analogue
circuitry in this book — a systematic way to solve any linear network
(node voltages), and a compression trick (Th\'evenin) that turns a whole
network into one source and one resistor. By the end you can take any
resistor circuit this book will build, write its equations by hand in a few
lines, and check the answer in the simulator before reaching for a wire.

\section{Series and parallel are about nodes}
\label{sec:seriesparallel}

\begin{definition}[Series]\index{series}
Two elements are in \emph{series} when they share exactly one node and nothing
else connects to that node. Whatever current enters one must leave through
the other — there is nowhere else for it to go.
\end{definition}

\begin{definition}[Parallel]\index{parallel}
Two elements are in \emph{parallel} when both of their terminals connect to the
same pair of nodes. They therefore see the same voltage.
\end{definition}

The definitions are about \emph{nodes}, not geometry. Figure~\ref{fig:redraws}
shows one circuit — $R_1$ in parallel with $R_2$ — drawn four ways. Beginners
call (c) and (d) "different circuits" because the resistors point differently
or sit in a triangle. They are the same circuit: in all four drawings every
$R_1$ terminal and every $R_2$ terminal lands on the same two nodes. The
schematic is the topology (Section~1.4); squint at connections, not shapes.

\begin{figure}[htbp]
\centering
% fig_redraws.tex -- Figure 3.1: same parallel network drawn four ways
\begin{tikzpicture}[scale=0.82, transform shape]
  % drawing 1
  \draw (0,4.5) to[short,o-] (0,5.3) -- (3,5.3) to[short,-o] (3,4.5);
  \draw (0.8,5.3) to[R=$R_1$] (0.8,3.5) -- (0.8,3.0);
  \draw (2.2,5.3) to[R=$R_2$] (2.2,3.5) -- (2.2,3.0);
  \draw (0.8,3.0) -- (2.2,3.0);
  \draw (0,3.0) to[short,o-] (0,4.5); \draw (3,3.0) to[short,-o] (3,4.5);
  \draw (0,3.0) -- (0.8,3.0); \draw (2.2,3.0) -- (3,3.0);
  \node[below] at (1.5,2.8) {(a)};
  % drawing 2: folded rail
  \draw (5.5,4.5) to[short,o-] (5.5,5.3) -- (8.5,5.3) to[short,-o] (8.5,4.5);
  \draw (6.3,5.3) to[R=$R_1$] (6.3,3.5) -- (6.3,3.0);
  \draw (7.7,5.3) -- (7.7,4.6) to[R=$R_2$] (9.1,4.6) -- (9.1,3.0) -- (6.3,3.0);
  \draw (5.5,3.0) to[short,o-] (5.5,4.5); \draw (8.5,3.0) to[short,-o] (8.5,4.5);
  \draw (5.5,3.0) -- (6.3,3.0); \draw (8.5,3.0) -- (8.5,3.0);
  \node[below] at (7.3,2.8) {(b)};
  % drawing 3: sideways
  \draw (11,4.5) to[short,o-] (11,5.3) -- (14,5.3) to[short,-o] (14,4.5);
  \draw (11.8,5.3) to[R=$R_1$] (13.4,5.3);
  \draw (11.8,3.6) to[R=$R_2$] (13.4,3.6);
  \draw (11.8,3.0) -- (11.8,5.3); \draw (13.4,3.0) -- (13.4,5.3);
  \draw (11.8,3.6) -- (13.4,3.6);
  \draw (11,3.0) to[short,o-] (11,4.5); \draw (14,3.0) to[short,-o] (14,4.5);
  \draw (11,3.0) -- (11.8,3.0); \draw (13.4,3.0) -- (14,3.0);
  \node[below] at (12.6,2.8) {(c)};
  % drawing 4: squashed triangle
  \draw (0.5,0.2) to[short,o-] (0.5,1.0) -- (1.5,1.8) -- (3.5,1.8) to[short,-o] (4.3,0.6);
  \draw (1.8,1.8) to[R=$R_1$] (1.8,0.0);
  \draw (3.2,1.8) to[R=$R_2$] (3.2,0.0);
  \draw (0.5,0.2) -- (1.8,0.0) -- (3.2,0.0) -- (4.3,0.6);
  \draw (4.3,0.6) to[short,-o] (4.3,1.0); \draw (4.3,1.0) -- (3.5,1.8);
  \node[below] at (2.4,-0.2) {(d)};
\end{tikzpicture}

\caption{One parallel circuit, four drawings. In every panel each end of
$R_1$ shares a node with an end of $R_2$ and nothing in between splits them.
If you can label only two distinct nodes, the resistors are in parallel —
however the wires are bent.}
\label{fig:redraws}
\end{figure}

\begin{exercise}[easy]
In Figure~\ref{fig:redraws}(b), trace from the top of $R_1$ to the top of
$R_2$ without lifting your pen. You never pass through a component — so those
ends are one node. Do the same for the bottoms.
\end{exercise}

\section{Equivalent resistance}
\label{sec:req}

A network of resistors between two terminals behaves, at those terminals, like
one resistor called the \emph{equivalent resistance}\index{equivalent
resistance} $R_{\mathrm{eq}}$. Two cases cover nearly everything.

\textbf{Series:} the same current crosses both, voltages add (KVL), so
$R_{\mathrm{eq}} = R_1 + R_2$. \textbf{Parallel:} the same voltage sits across
both, currents add (KCL):
\[
\frac{1}{R_{\mathrm{eq}}} = \frac{1}{R_1} + \frac{1}{R_2}
\qquad\Longleftrightarrow\qquad
R_{\mathrm{eq}} = \frac{R_1 R_2}{R_1 + R_2}\quad\text{("product over sum")}.
\]
For $n$ resistors, series just sums; parallel sums the reciprocals. Reduce
any network by repeatedly collapsing the two-terminal pieces. Sanity checks
that catch almost every slip: in parallel, $R_{\mathrm{eq}}$ is
\emph{smaller than the smallest} resistor (more paths, more current); in
series, $R_{\mathrm{eq}}$ is larger than the largest.

\begin{exercise}[easy]
Reduce this to one resistor: \qty{1}{\kilo\ohm} in parallel with
\qty{2}{\kilo\ohm}, in series with \qty{470}{\ohm}. (Answer:
$667 + 470 \approx \qty{1.14}{\kilo\ohm}$.)
\end{exercise}

\section{The voltage divider}
\label{sec:divider}

Two series resistors across a source split the voltage in proportion to their
resistances. With $R_1$ on top and $R_2$ to ground as in
Figure~\ref{fig:divider},
\[
v_{\mathrm{out}} \;=\; V_{\mathrm{in}}\,\frac{R_2}{R_1+R_2}.
\]
One line of KVL derives it: the current is $V_{\mathrm{in}}/(R_1+R_2)$, and
the output is that current times $R_2$. This fragment appears dozens of times
in this book — scaling sensor signals, setting logic levels, biasing
transistors. Learn it once, recognise it everywhere.

\begin{figure}[htbp]
\centering
% fig_divider.tex -- Figure 3.2: annotated voltage divider
\begin{circuitikz}[american]
  \draw (0,0) to[battery1, l_=$V_{\mathrm{in}}$, invert] (0,4);
  \draw (0,4) to[R, l=$R_1$, v^>=$v_{R1}$, i>_=$i$] (3,4);
  \draw (3,4) to[R, l=$R_2$, v^>=$v_{\mathrm{out}}$, i>_=$i$] (3,0) -- (0,0);
  \draw (3,4) to[short,-o] (4.3,4) node[right]{$v_{\mathrm{out}}$};
  \node[right,align=left] at (4.8,2.6) {$v_{\mathrm{out}} = V_{\mathrm{in}}\,\dfrac{R_2}{R_1+R_2}$\\[4pt] $i = \dfrac{V_{\mathrm{in}}}{R_1+R_2}$};
\end{circuitikz}

\caption{The voltage divider. The output is taken \emph{at the middle node};
$R_2$ is the resistor \emph{to ground}. Notice the formula needs the same
current through both — the moment something draws current out of the middle
node, $v_{\mathrm{out}}$ falls (Section~\ref{sec:thevenin} quantifies it).}
\label{fig:divider}
\end{figure}

\begin{example}[A 3.3 V logic level from a 5 V rail]
\label{ex:div33}
\textbf{Strategy.} The ratio sets the voltage; the absolute values set how
much current the divider wastes. Pick a ratio from the E24 ladder, then
choose a magnitude that burns little current but still makes the output
stiff compared to anything the logic input draws (chip inputs want microamps
or less).

\textbf{Solution.} We need $R_2/(R_1+R_2) = 3.3/5 = 0.66$, i.e.
$R_1 : R_2 \approx 1 : 1.94$. E24 gives a clean hit: $R_1=\qty{1}{\kilo\ohm}$,
$R_2=\qty{2}{\kilo\ohm}$, ratio $2000/3000=0.667$.
Then $v_{\mathrm{out}}=5\times0.667=\qty{3.33}{\volt}$ and the standing
current is $5/3000=\qty{1.67}{\milli\ampere}$ — under \qty{10}{\milli\watt}
from the rail. Tolerance corners (\qty{5}{\percent} parts): worst high is
$5\times2.1/3.05=\qty{3.44}{\volt}$, worst low is
$5\times1.9/2.95=\qty{3.22}{\volt}$ — a \qty{3.3}{\volt} input typically
accepts down to about \qty{2.5}{\volt} as "high", so both corners pass.

\textbf{Check.} Listing~\ref{lst:dividerop} runs it. Notice the simulator
takes the question "what are all node voltages and branch currents" literally
— and answers 3.33 V and $1.67$ mA exactly as derived.
\end{example}

\begin{lstlisting}[language={},caption={\texttt{code/ch03\_divider.cir} — the divider as a netlist.},label=lst:dividercir]
* ch03_divider.cir -- 5 V -> 3.3 V divider, operating point
V1 1 0 DC 5
R1 1 2 1000
R2 2 0 2000
.op
.end
\end{lstlisting}

\begin{lstlisting}[language={},caption={ngspice output of \texttt{code/ch03\_divider.cir}, verbatim (simulation S1).},label=lst:dividerop]
V(2)                             3.333333e+00
V(1)                             5.000000e+00
v1#branch                        -1.66667e-03
\end{lstlisting}

\section{The current divider}
\label{sec:curdiv}

Parallel is the mirror image. Two parallel resistors share a voltage and split
an incoming current \emph{inversely} to their resistances — more current takes
the easier path:
\[
i_2 = I\,\frac{R_1}{R_1+R_2}.
\]
Note the twist: the current through $R_2$ is proportional to the
\emph{other} resistor. (The clean derivation uses conductances
$G=1/R$\index{conductance}: in parallel, conductances add and current splits
in proportion to $G$, i.e. $i_2 = I\,G_2/(G_1+G_2)$ — same thing, one less
fraction to fumble.)

\section{KCL and KVL: bookkeeping, signed}
\label{sec:klaws}

\begin{definition}[Kirchhoff's current law]\index{Kirchhoff's current law}
The sum of currents \emph{into} a node is zero — or equivalently, what flows
in flows out. Charge is conserved and nodes store none of it.
\end{definition}

\begin{definition}[Kirchhoff's voltage law]\index{Kirchhoff's voltage law}
The sum of signed voltage changes around any closed loop is zero. Voltage is
a potential difference; returning to the same node must return the same
potential.
\end{definition}

The laws are trivial. The discipline is the \emph{signs}. Pick a direction and
label every drop before writing equations — Figure~\ref{fig:kvlkcl} marks the
convention this book uses: a drop is $+$ when your loop direction enters the
$+$ labelled terminal. Then the equation writes itself. Sign errors, not
concept errors, are what actually go wrong (Section~\ref{sec:mistakes3}).

\begin{figure}[htbp]
\centering
% fig_kvlkcl.tex -- Figure 3.3: KVL signed drops + KCL node
\begin{circuitikz}[american]
  % left: KVL loop
  \draw (0,0) to[battery1, l_=\qty{9}{\volt}, invert] (0,3.4);
  \draw (0,3.4) to[R, l=\qty{220}{\ohm}, v^<=$v_1$] (3,3.4);
  \draw (3,3.4) to[R, l=\qty{470}{\ohm}, v^<=$v_2$] (3,0) -- (0,0);
  \draw[->,thick] (0.9,1.2) arc[start angle=200,end angle=-20,radius=0.75];
  \node at (1.55,1.5) {$\sum v_k = 0$};
  % right: KCL node
  \node[circ] (n) at (7,2.2) {};
  \draw (n) -- (5.4,3.4) node[left]{$i_1$};
  \draw (n) -- (8.6,3.4) node[right]{$i_2$};
  \draw (n) -- (7,0.8) node[below]{$i_3$};
  \draw[->,thick] (5.9,3.2) -- (6.8,2.4);
  \draw[->,thick] (8.1,3.2) -- (7.2,2.4);
  \draw[->,thick] (7,1.6) -- (7,2.0);
  \node[right,align=left] at (8.8,1.6) {signs first:\\ $+i_1+i_2-i_3=0$\\ if $i_3$ leaves};
  \node[below] at (9.6,0.6) {};
\end{circuitikz}

\caption{Left: KVL on the series loop — walking clockwise,
$+9 - v_1 - v_2 = 0$. Right: KCL at a node — assign each current a direction
\emph{first} (here $i_3$ defined leaving), then sum algebraically:
$i_1 + i_2 - i_3 = 0$. Negative answers just mean "the arrow was guessed
backwards", which is fine.}
\label{fig:kvlkcl}
\end{figure}

\begin{checkpoint}
In Figure~\ref{fig:kvlkcl} (left), $v_1=\qty{3}{\volt}$. Write the KVL sum
and find $v_2$ without computing any current. (Answer:
$9-3-v_2=0 \Rightarrow v_2=\qty{6}{\volt}$.)
\end{checkpoint}

\section{The node-voltage method}
\label{sec:nodevoltage}

For a network that will not collapse to one $R_{\mathrm{eq}}$, do this
instead\index{node-voltage method}: pick a ground, name every other node
voltage $v_1, v_2, \dots$, and write KCL at each node with each current
expressed via Ohm's law. One equation per node, solved together — that is the
entire method, and it is exactly what the simulator does internally.

\begin{example}[Two sources, one node: by hand then by ngspice]
\label{ex:twoloop}
\textbf{Strategy.} Node 3 joins three resistors to two sources. Write KCL
there; the branch currents are $(v_3 - 12)/100$, $(v_3 - 6)/220$ and
$v_3/470$. One unknown, one equation.

\textbf{Solution.}
\[
\frac{v_3-12}{100} + \frac{v_3-6}{220} + \frac{v_3}{470} = 0
\;\Longrightarrow\;
v_3\!\left(\tfrac{1}{100}+\tfrac{1}{220}+\tfrac{1}{470}\right)
= \tfrac{12}{100}+\tfrac{6}{220}
\]
\[
v_3 = \frac{0.1473}{0.01667} \approx \qty{8.83}{\volt}.
\]
Notice the form: \emph{weighted average}. $v_3$ lies between 6 and 12 V,
pulled toward the 12 V source because its \qty{100}{\ohm} link is the
strongest (lowest resistance).

\textbf{Check.} Listing~\ref{lst:twoloopop} — ngspice agrees to the fourth
decimal, and reports the branch currents $-31.7$ mA and $+12.9$ mA that the
hand work predicts ($\frac{8.83-12}{100}=-31.7$ mA;
$\frac{8.83-6}{220}=+12.9$ mA; both plus $v_3/470=18.8$ mA sum to zero —
KCL, closing the loop on itself).
\end{example}

\begin{lstlisting}[language={},caption={\texttt{code/ch03\_twoloop.cir} — the two-source network of Example~\ref{ex:twoloop}.},label=lst:twoloopcir]
* ch03_twoloop.cir -- two sources feeding one node through resistors
V1 1 0 DC 12
V2 2 0 DC 6
R1 1 3 100
R2 2 3 220
R3 3 0 470
.op
.end
\end{lstlisting}

\begin{lstlisting}[language={},caption={ngspice output of \texttt{code/ch03\_twoloop.cir}, verbatim (simulation S2).},label=lst:twoloopop]
V(3)                             8.832947e+00
V(2)                             6.000000e+00
V(1)                             1.200000e+01
v1#branch                        -3.16705e-02
v2#branch                        1.287703e-02
\end{lstlisting}

\section{Th\'evenin and Norton: networks compressed}
\label{sec:thevenin}

\begin{definition}[Th\'evenin equivalent]\index{Th\'evenin equivalent}
Any linear network viewed from two terminals behaves exactly like one voltage
source $V_{\mathrm{th}}$ in series with one resistor $R_{\mathrm{th}}$:
$V_{\mathrm{th}}$ is the open-circuit voltage at the terminals, and
$R_{\mathrm{th}}$ is the resistance looking in (all independent sources
zeroed) — equivalently $R_{\mathrm{th}} = V_{\mathrm{th}}/I_{\mathrm{sc}}$
where $I_{\mathrm{sc}}$ is the short-circuit current.
\end{definition}

"Behaves exactly like" means \emph{for every possible load}: the terminal's
voltage-versus-load-current line is identical. Figure~\ref{fig:thevenin}
shows the picture; Figure~\ref{fig:loadline} proves it on the divider —
sweep a load current through the original and through the model and the lines
coincide. The \emph{Norton equivalent}\index{Norton equivalent} is the same
idea drawn the other way: a current source $I_{\mathrm{sc}}$ in parallel with
$R_{\mathrm{th}}$, interchangeable by $V_{\mathrm{th}}=I_{\mathrm{sc}}R_{\mathrm{th}}$.
Worked use of it: our divider's Norton form is \qty{5}{\milli\ampere} in
parallel with \qty{667}{\ohm}. Ask how much current a \qty{2}{\kilo\ohm} load
draws and the current divider answers in one line:
$5\,\mathrm{mA}\times 667/(667+2000) = \qty{1.25}{\milli\ampere}$ — the same
answer KCL gives on the full circuit ($v_{\mathrm{out}}=\qty{2.5}{\volt}$),
with the network replaced by two numbers. Norton is the natural form
whenever the quantity you care about is current — sensors, current-output
devices, parallel paths.

\begin{figure}[htbp]
\centering
% fig_thevenin.tex -- Figure 3.4: Thevenin equivalent picture
\begin{tikzpicture}
  % original network
  \draw[draw=gray!60,fill=gray!8,rounded corners] (-0.2,-0.4) rectangle (2.9,2.4);
  \node at (1.35,2.1) {\small any linear network};
  \draw (0.4,0) -- (0.4,1.6) node[circ,label=above:$a$]{};
  \draw (2.3,0) -- (2.3,1.6) node[circ,label=above:$b$]{};
  % arrow
  \draw[->,thick] (3.2,1.0) -- (4.3,1.0) node[midway,above]{\small measure};
  % thevenin
  \draw (5.0,1.6) to[battery1, l_=$V_{\mathrm{th}}$, invert] (5.0,0);
  \draw (5.0,1.6) to[R, l=$R_{\mathrm{th}}$] (6.9,1.6) node[circ,label=above:$a$]{};
  \draw (6.9,0) node[circ,label=above:$b$]{} -- (5.0,0);
  \node[below,align=center] at (6.0,-0.6) {$V_{\mathrm{th}}$ = open-circuit voltage\\ $R_{\mathrm{th}}$ = resistance looking in};
\end{tikzpicture}

\caption{The Th\'evenin picture: whatever is inside the left box — any number
of sources and resistors — the terminals see only $V_{\mathrm{th}}$ and
$R_{\mathrm{th}}$. Two measurements determine it: open-circuit voltage, and
either short-circuit current or resistance-with-sources-zeroed.}
\label{fig:thevenin}
\end{figure}

\begin{example}[The divider, compressed — and when "3.3 V" lies]
\label{ex:thevdiv}
\textbf{Strategy.} Compress the Example~\ref{ex:div33} divider
($R_1=\qty{1}{\kilo\ohm}$, $R_2=\qty{2}{\kilo\ohm}$ on 5 V) into a
Th\'evenin model, then ask it a loaded question.

\textbf{Solution.} $V_{\mathrm{th}}$ = open-circuit output = \qty{3.33}{\volt}.
Looking in with the 5 V source zeroed (shorted): $R_1$ and $R_2$ are both
from the output node to ground $\Rightarrow
R_{\mathrm{th}} = 1\|2\,\mathrm{k}\Omega = \qty{667}{\ohm}$.
Now attach a \qty{1}{\kilo\ohm} load: the model predicts
$v_{\mathrm{out}} = 3.33 \times 1000/1667 = \qty{2.0}{\volt}$ — the "3.3 V
rail" collapsed by \qty{40}{\percent}. The simplification does not lie; the
lie was treating an unstiff divider as a voltage source. The load-line
sweep in Figure~\ref{fig:loadline} (from the Listings below) shows the model
and the original agreeing to seven digits: slope $-667\ \Omega$, intercept
\qty{3.33}{\volt}, zero crossing at $I_{\mathrm{sc}}=\qty{5}{\milli\ampere}
= 3.33/0.667$.

\textbf{Rule of thumb.} A divider feeding a load wants
$R_{\mathrm{th}} \ll R_{\mathrm{load}}$ — a factor of 10 keeps the sag under
about \qty{9}{\percent}.
\end{example}

\begin{lstlisting}[language={},caption={\texttt{code/ch03\_thevenin.cir} — sweep a load current, record $v_{\mathrm{out}}$ (simulation S3). \texttt{code/ch03\_thevenin\_model.cir} is the same file with the two resistors replaced by $V_{\mathrm{th}}{=}\qty{3.333}{\volt}$ and $R_{\mathrm{th}}{=}\qty{667}{\ohm}$.},label=lst:thevcir]
* ch03_thevenin.cir -- load-line sweep: original divider
I1 2 0 DC 0          * load current pulled OUT of node 2
V1 1 0 DC 5
R1 1 2 1000
R2 2 0 2000
.control
dc i1 0 0.01 0.0005
let vout = v(2)
set filetype=ascii
wrdata ch03_divider_loadline.dat vout
.endc
.end
\end{lstlisting}

\begin{lstlisting}[language={},caption={Verification output, verbatim — every 4th row shown of 21 simulated rows: load current, $v_{\mathrm{out}}$ of the real divider, $v_{\mathrm{out}}$ of the Th\'evenin model. Column 3 minus column 2 never exceeds \qty{7e-8}{\volt} on any row.},label=lst:thevout]
i_load          v_orig          v_model
0.00000000e+00  3.33333333e+00  3.33333330e+00
2.00000000e-03  2.00000000e+00  1.99999996e+00
4.00000000e-03  6.66666667e-01  6.66666620e-01
6.00000000e-03 -6.66666667e-01 -6.66666720e-01
8.00000000e-03 -2.00000000e+00 -2.00000006e+00
1.00000000e-02 -3.33333333e+00 -3.33333340e+00
\end{lstlisting}

\begin{figure}[htbp]
\centering
% fig_loadline.tex -- Figure 3.6: load-line sweep, divider vs Thevenin model
\begin{tikzpicture}
\begin{axis}[width=0.85\textwidth, height=6.2cm,
  xlabel={load current (\unit{\milli\ampere})},
  ylabel={$v_{\mathrm{out}}$ (\unit{\volt})},
  xmin=0, xmax=10, ymin=0, ymax=3.6,
  grid=both, grid style={gray!15},
  legend style={at={(0.97,0.30)},anchor=east,font=\small}]
  \addplot[blue, thick] table[x index=0,y index=1,y expr={\thisrowno{0}*1}] {figures/ch03/ch03_loadline.dat};
  \addlegendentry{original divider}
  \addplot[red, dashed, thick] table[x index=0,y index=1] {figures/ch03/ch03_loadline_model.dat};
  \addlegendentry{Th\'evenin model}
  \addplot[black, only marks, mark=x, mark size=3pt] coordinates {(0,3.333) (5,0)};
  \node[font=\small, anchor=west] at (axis cs:0.25,3.42) {$V_{\mathrm{th}}=\qty{3.33}{\volt}$};
  \node[font=\small, anchor=south west] at (axis cs:5.15,0.02) {$I_{\mathrm{sc}}=\qty{5.0}{\milli\ampere}$};
\end{axis}
\end{tikzpicture}

\caption{Simulation S3: $v_{\mathrm{out}}$ against pulled load current for the
real divider (blue) and its Th\'evenin model (red, dashed) — the same line,
to \qty{70}{\nano\volt}. The marked endpoints are the theorem's two handles:
open-circuit voltage \qty{3.33}{\volt} and short-circuit current
\qty{5}{\milli\ampere}. Every linear network has exactly one such line.}
\label{fig:loadline}
\end{figure}

\section{Measuring without breaking the circuit}
\label{sec:meterloading}

A voltmeter connects in parallel — so it adds a resistance in parallel.
Real DMMs present about \qty{10}{\mega\ohm}\index{meter loading}; on
kiloohm-scale circuits that is invisible, but measure a megaohm divider and
the meter \emph{becomes part of the circuit}
(Figure~\ref{fig:meter}): $1\,\mathrm{M}\Omega \| 10\,\mathrm{M}\Omega =
\qty{909}{\kilo\ohm}$, so a \qty{5}{\volt} rail reads $5\times909/1909 =
\qty{2.38}{\volt}$ instead of \qty{2.5}{\volt}. The meter did not "use up"
voltage; it changed the divider. When a reading surprises you, suspect the
probe before the circuit.

\begin{figure}[htbp]
\centering
% fig_meter.tex -- Figure 3.5: meter loading
\begin{circuitikz}[american]
  \draw (0,0) to[battery1, l_=\qty{5}{\volt}, invert] (0,3.4);
  \draw (0,3.4) to[R, l=\qty{1}{\mega\ohm}] (3,3.4);
  \draw (3,3.4) to[R, l_=\qty{1}{\mega\ohm}] (3,0) -- (0,0);
  % DMM across bottom resistor
  \draw (3,3.4) -- (5,3.4) to[vR, l_=DMM, n=mymeter] (5,0) -- (3,0);
  \draw[dashed] (5,3.4) to[R, l=\qty{10}{\mega\ohm}] (5,0);
  \node[right,align=left] at (5.9,1.9) {reads \qty{2.38}{\volt},\\ not \qty{2.5}{\volt}};
\end{circuitikz}

\caption{Meter loading. Two \qty{1}{\mega\ohm} resistors divide 5 V to
\qty{2.5}{\volt} until a \qty{10}{\mega\ohm} meter joins the bottom leg —
\qty{2.38}{\volt} now. Dashed: the resistance you added by measuring.}
\label{fig:meter}
\end{figure}

\begin{exercise}[medium — predict, then run]
A \qty{4.7}{\mega\ohm} / \qty{4.7}{\mega\ohm} divider sits on \qty{9}{\volt}.
Predict the DMM reading ($R_{\mathrm{meter}}=\qty{10}{\mega\ohm}$), write the
three-line netlist, run it, compare. (Answer: $4.7\|10=\qty{3.20}{\mega\ohm}$,
$9\times3.20/7.90=\qty{3.64}{\volt}$ — the meter drags it down by
\qty{19}{\percent}.)
\end{exercise}

\section{Common mistakes}
\label{sec:mistakes3}

\begin{tripwire}[KVL around a non-loop / KCL at a non-node]
Writing $\sum v = 0$ across elements that do not form a closed loop, or
summing currents at a "node" that secretly contains a junction you missed.
Trace the loop with a finger; list the node's wires explicitly.
\end{tripwire}

\begin{tripwire}[Signs guessed mid-equation]
Assigning polarities while solving instead of before. If $v_1$ was labelled
one way in the figure and flipped in the algebra, the sign error compounds
silently. Label first; solve second; never adjust.
\end{tripwire}

\begin{tripwire}[Using the divider formula on a loaded divider]
$v_{\mathrm{out}} = V R_2/(R_1+R_2)$ assumes the middle node's only currents
are through $R_1$ and $R_2$. Any load at the node breaks the premise —
Example~\ref{ex:thevdiv}'s \qty{2.0}{\volt} instead of \qty{3.3}{\volt}.
Check what else connects.
\end{tripwire}

\begin{tripwire}[Mixing series and parallel formulas mid-redraw]
$R_1+R_2$ used where resistors share a node-pair but also a third element,
or "product over sum" applied to resistors that are not both on the same two
nodes. Redraw to the canonical form of Figure~\ref{fig:redraws} before any
arithmetic.
\end{tripwire}

\section{Chapter summary}
\label{sec:ch3summary}

Figure~\ref{fig:cheatsheet} is the whole chapter on one card. The method
hierarchy: redraw until it is series/parallel; if it will not collapse, write
KCL at each node; if you only care about one port, compress to Th\'evenin.
Everything else is algebra.

\begin{figure}[htbp]
\centering
% fig_cheatsheet.tex -- Figure 3.7: chapter cheat-sheet table (used as summary)
\begin{tabular}{@{}llll@{}}
\toprule
Quantity & Series & Parallel & Note \\
\midrule
Resistance & $R_{\mathrm{eq}}=R_1+R_2$ & $R_{\mathrm{eq}}=\dfrac{R_1R_2}{R_1+R_2}$ & two at a time \\
Voltage & splits $\propto R$ & same on each & KVL \\
Current & same through each & splits $\propto 1/R$ & KCL \\
Divider & $v_{\mathrm{out}}=V\dfrac{R_2}{R_1+R_2}$ & $i_2=I\dfrac{R_1}{R_1+R_2}$ & unloaded! \\
Th\'evenin & $V_{\mathrm{th}}=v_{\mathrm{oc}}$ & $R_{\mathrm{th}}=V_{\mathrm{th}}/I_{\mathrm{sc}}$ & one line fits all loads \\
\bottomrule
\end{tabular}

\caption{Cheat-sheet for Chapter 3. The last row is why Th\'evenin matters:
one measured line replaces the whole internal network.}
\label{fig:cheatsheet}
\end{figure}

\begin{wins}
\begin{itemize}
\item You classify connections by node counting and reduce any
series/parallel network to one resistor.
\item You can design and load-check a voltage divider — the most common two
resistors in this book.
\item KCL, KVL and the node-voltage method solve any linear resistive
circuit; you verified one against ngspice to the fourth decimal.
\item Th\'evenin compression tells you when a "supply" will sag, and the
meter-loading rule keeps your measurements honest.
\end{itemize}
\end{wins}

\section*{Exercises}
\addcontentsline{toc}{section}{Exercises}

\begin{exercise}[easy]
Three resistors — \qty{1}{\kilo\ohm}, \qty{2}{\kilo\ohm}, \qty{3}{\kilo\ohm} —
are placed end-to-end in a ring, terminals taken across the
\qty{1}{\kilo\ohm} only. What is $R_{\mathrm{eq}}$? (Careful: it is a
parallel question.)
\end{exercise}

\begin{exercise}[easy]
A node has four wires carrying \qty{2}{\milli\ampere},
\qty{-5}{\milli\ampere} and \qty{4}{\milli\ampere} (signs = into node).
What does the fourth carry?
\end{exercise}

\begin{exercise}[easy]
Using Figure~\ref{fig:divider}'s labels: $V_{\mathrm{in}}=\qty{12}{\volt}$,
$R_1=\qty{4.7}{\kilo\ohm}$, $R_2=\qty{10}{\kilo\ohm}$. Find
$v_{\mathrm{out}}$ and the standing current.
\end{exercise}

\begin{exercise}[easy]
True or false, with a one-sentence justification: two \qty{100}{\ohm}
resistors in parallel carry the same current as one \qty{50}{\ohm} resistor
at the same voltage.
\end{exercise}

\begin{exercise}[medium — derive]
Derive the parallel law from KCL: for $n$ resistors on the same node pair at
voltage $v$, show $1/R_{\mathrm{eq}} = \sum 1/R_k$; specialise to the
product-over-sum two-resistor form.
\end{exercise}

\begin{exercise}[medium — the classic trap]
Design a divider that outputs \qty{3.3}{\volt} from \qty{5}{\volt}
\emph{while the output node sources \qty{10}{\milli\ampere}}. Hint: you
cannot treat the load as infinite — write KCL at the output node including
the \qty{10}{\milli\ampere}. Choose real E24 values and report the delivered
voltage.
\end{exercise}

\begin{exercise}[medium]
Find the Th\'evenin equivalent of: \qty{9}{\volt} source, \qty{2}{\kilo\ohm}
series, \qty{3}{\kilo\ohm} to ground, terminals at the middle node. Then
predict the terminal voltage with a \qty{6}{\kilo\ohm} load.
\end{exercise}

\begin{exercise}[medium]
You probe the middle of an unloaded \qty{100}{\kilo\ohm}/\qty{100}{\kilo\ohm}
divider on \qty{10}{\volt} and read \qty{4.762}{\volt}. Infer the meter's
input resistance.
\end{exercise}

\begin{exercise}[hard — mesh]
The circuit of Example~\ref{ex:twoloop} modified: add a fourth resistor
\qty{330}{\ohm} from node 3 to a \emph{new} node 4, and connect node 4 to
ground through \qty{680}{\ohm}. Write the two node-voltage equations, solve
by hand, then confirm with a netlist you write yourself.
\end{exercise}

\begin{exercise}[hard]
A linear network is characterised only by two measurements: open-circuit
\qty{4.2}{\volt}; with a \qty{1}{\kilo\ohm} load it reads \qty{3.0}{\volt}.
Reconstruct $V_{\mathrm{th}}$ and $R_{\mathrm{th}}$, and state what load
resistance would pull the terminal down to half its open-circuit value.
\end{exercise}

\section*{Hints and answers}
\addcontentsline{toc}{section}{Hints and answers}

\begin{description}
\item[Ring.] \qty{1}{\kilo\ohm} in parallel with the series chain
$(2+3)\,\mathrm{k}\Omega$: $R_{\mathrm{eq}}=1000\cdot5000/6000
\approx\qty{833}{\ohm}$.
\item[KCL.] The three given sum to $+1$ mA; the fourth must carry
\qty{-1}{\milli\ampere} (i.e.\ \qty{1}{\milli\ampere} leaving).
\item[12 V divider.] $v_{\mathrm{out}}=12\times10/14.7=\qty{8.16}{\volt}$;
$i=12/14.7\,\mathrm{k}=\qty{0.816}{\milli\ampere}$.
\item[50 ohm.] True — equivalence means identical terminal behaviour; that is
the definition's whole point.
\item[Derivation.] Same $v$ on all $n$ branches, $i_k=v/R_k$, KCL gives
$I=v\sum 1/R_k$; compare with $I=v/R_{\mathrm{eq}}$. For $n=2$:
$R_{\mathrm{eq}}=1/(1/R_1+1/R_2)=R_1R_2/(R_1+R_2)$.
\item[Loaded divider.] KCL at the output: $(5-v)/R_1 = v/R_2 + 0.01$ with
$v=3.3$. Pick $R_1=\qty{100}{\ohm}$: $17$ mA in, so $3.3/R_2=7$ mA,
$R_2\approx\qty{471}{\ohm}\to\qty{470}{\ohm}$ E24; simulated value
$\approx\qty{3.30}{\volt}$. Any pair with the same current budget works.
\item[Th\'evenin.] $V_{\mathrm{th}}=9\times3/5=\qty{5.4}{\volt}$;
$R_{\mathrm{th}}=2\|3=\qty{1.2}{\kilo\ohm}$; loaded:
$5.4\times6/7.2=\qty{4.5}{\volt}$.
\item[Meter inference.] Reading $4.762=10\times R_{\mathrm{eq}}/(100+R_{\mathrm{eq}})$
gives $R_{\mathrm{eq}}\approx\qty{90.9}{\kilo\ohm}$, so
$R_{\mathrm{meter}}\approx\qty{1}{\mega\ohm}$.
\item[Mesh/nodes.] Equations: $(v_3-12)/100+(v_3-6)/220+v_3/470+(v_3-v_4)/330=0$;
$v_4/680+(v_4-v_3)/330=0$. Second gives $v_4=0.673\,v_3$; substituting,
$v_3\approx\qty{8.34}{\volt}$, $v_4\approx\qty{5.61}{\volt}$.
\item[Two measurements.] $V_{\mathrm{th}}=\qty{4.2}{\volt}$; under
\qty{1}{\kilo\ohm}: $3.0=4.2\times1000/(R_{\mathrm{th}}+1000)$ gives
$R_{\mathrm{th}}=\qty{400}{\ohm}$; half open-circuit when
$R_{\mathrm{load}}=R_{\mathrm{th}}=\qty{400}{\ohm}$.
\end{description}

\chapternotes
\begin{enumerate}
\item Series/parallel-by-nodes presentation and the four-redraws exercise
format follow the sequencing of Kuphaldt's Vol.~I ch.~5--6 \cite{aacvol1};
the example networks are original constructions.
\item The node-voltage method and the weighted-average reading of
Example~\ref{ex:twoloop} track MIT 6.002 lectures 1--3 \cite{ocw6002}, which
also teach Th\'evenin as a load-line measurement — the basis of
Figure~\ref{fig:loadline}'s simulation.
\item Meter-loading numbers use the conventional \qty{10}{\mega\ohm} DMM
input resistance \cite{sparkfunmultimeter}.
\item All ngspice outputs in this chapter are verbatim from ngspice-47 on the
author's machine; the load-line agreement ($<\qty{1e-7}{\volt}$) is exact up
to the model's truncated $V_{\mathrm{th}}$ \cite{ngspice}.
\end{enumerate}
